\documentclass[a4paper,12pt]{article}

% Поддержка русского языка и кодировок
\usepackage[utf8]{inputenc}
\usepackage[T2A]{fontenc}
\usepackage[russian]{babel}

% Математические пакеты
\usepackage{amsmath, amsfonts, amssymb, amsthm}

% Настройка полей
\usepackage[left=3cm, right=1.5cm, top=2cm, bottom=2cm]{geometry}

% Пакеты для графики и таблиц
\usepackage{graphicx}
\usepackage{tabularx}
\usepackage{booktabs}
\usepackage{float}

% Пакет для вставки исходного кода
\usepackage{listings}
\usepackage{xcolor}

\lstset{
    language=C++,
    basicstyle=\ttfamily\footnotesize,
    keywordstyle=\color{blue}\bfseries,
    stringstyle=\color{red},
    commentstyle=\color{gray}\itshape,
    numbers=left,
    numberstyle=\tiny,
    stepnumber=1,
    numbersep=5pt,
    backgroundcolor=\color{white},
    showspaces=false,
    showstringspaces=false,
    showtabs=false,
    frame=single,
    tabsize=4,
    captionpos=b,
    breaklines=true,
    breakatwhitespace=false,
    escapeinside={\%*}{*)},
    extendedchars=true,
    keepspaces=true,
    literate={а}{{\cyra}}1 {б}{{\cyrb}}1 {в}{{\cyrv}}1 {г}{{\cyrg}}1 {д}{{\cyrd}}1 
             {е}{{\cyre}}1 {ё}{{\"e}}1 {ж}{{\cyrzh}}1 {з}{{\cyrz}}1 {и}{{\cyri}}1 
             {й}{{\cyrishrt}}1 {к}{{\cyrk}}1 {л}{{\cyrl}}1 {м}{{\cyrm}}1 {н}{{\cyrn}}1 
             {о}{{\cyro}}1 {п}{{\cyrp}}1 {р}{{\cyrr}}1 {с}{{\cyrs}}1 {т}{{\cyrt}}1 
             {у}{{\cyru}}1 {ф}{{\cyrf}}1 {х}{{\cyrh}}1 {ц}{{\cyrc}}1 {ч}{{\cyrch}}1 
             {ш}{{\cyrsh}}1 {щ}{{\cyrshch}}1 {ъ}{{\cyrhrdsn}}1 {ы}{{\cyrery}}1 {ь}{{\cyrsftsn}}1 
             {э}{{\cyreve}}1 {ю}{{\cyryu}}1 {я}{{\cyrya}}1 
             {А}{{\CYRA}}1 {Б}{{\CYRB}}1 {В}{{\CYRV}}1 {Г}{{\CYRG}}1 {Д}{{\CYRD}}1 
             {Е}{{\CYRE}}1 {Ё}{{\"E}}1 {Ж}{{\CYRZH}}1 {З}{{\CYRZ}}1 {И}{{\CYRI}}1 
             {Й}{{\CYRISHRT}}1 {К}{{\CYRK}}1 {Л}{{\CYRL}}1 {М}{{\CYRM}}1 {Н}{{\CYRN}}1 
             {О}{{\CYRO}}1 {П}{{\CYRP}}1 {Р}{{\CYRR}}1 {С}{{\CYRS}}1 {Т}{{\CYRT}}1 
             {У}{{\CYRU}}1 {Ф}{{\CYRF}}1 {Х}{{\CYRH}}1 {Ц}{{\CYRC}}1 {Ч}{{\CYRCH}}1 
             {Ш}{{\CYRSH}}1 {Щ}{{\CYRSHCH}}1 {Ъ}{{\CYRHRDSN}}1 {Ы}{{\CYRERY}}1 {Ь}{{\CYRSFTSN}}1 
             {Э}{{\CYREVE}}1 {Ю}{{\CYRYU}}1 {Я}{{\CYRYA}}1
}

\begin{document}

\setcounter{page}{2}

% --- ОСНОВНОЙ ТЕКСТ ---
\section{Цель работы}
Разработать программу решения двумерной эллиптической краевой задачи методом конечных разностей. Протестировать программу на полиномах и численно оценить порядок аппроксимации. Область имеет Т-образную форму. Предусмотреть учет краевых условий первого и третьего рода.

\section{Теоретическая часть}
Метод конечных разностей основан на замене непрерывной расчетной области дискретной сеткой узлов и аппроксимации дифференциальных операторов разностными отношениями с использованием разложения функций в ряд Тейлора.

Для двумерного уравнения эллиптического типа:
\begin{equation}
-\lambda \Delta u + \gamma u = f(x, y)
\end{equation}

Дискретный аналог на неравномерной прямоугольной сетке определяется пятиточечным разностным выражением:
\begin{equation}
\begin{split}
\Delta_h u_{i,j} = \frac{2u_{i-1,j}}{h_{i-1}^x(h_i^x+h_{i-1}^x)} + \frac{2u_{i+1,j}}{h_i^x(h_i^x+h_{i-1}^x)} + \frac{2u_{i,j-1}}{h_{j-1}^y(h_j^y+h_{j-1}^y)} + \frac{2u_{i,j+1}}{h_j^y(h_j^y+h_{j-1}^y)} - \\ - \left(\frac{2}{h_{i-1}^x h_i^x} + \frac{2}{h_{j-1}^y h_j^y}\right)u_{i,j}
\end{split}
\end{equation}

Учет краевых условий осуществляется следующим образом:
Для узлов на границе $S_1$, где заданы условия первого рода $u = u_g$, диагональные элементы матрицы заменяются на 1, а элемент вектора правой части заменяется на значение $u_g$.
Для условия третьего рода $\lambda \frac{\partial u}{\partial n} + \beta(u - u_\beta) = 0$ производная по нормали аппроксимируется односторонней разностью первого порядка. Например, для верхней границы (внешняя нормаль сонаправлена с осью OY):
\begin{equation}
\lambda \frac{u_{i,j} - u_{i,j-1}}{h_{j-1}^y} + \beta u_{i,j} = \beta u_\beta \implies \left(\frac{\lambda}{h_{j-1}^y} + \beta\right)u_{i,j} - \frac{\lambda}{h_{j-1}^y}u_{i,j-1} = \beta u_\beta
\end{equation}

\section{Практическая часть}

\subsection{Построение сетки в Т-образной области}
Для решения задачи используется регулярная прямоугольная сетка, построенная в минимальном ограничивающем прямоугольнике $\Omega_{rect}$. 
Т-образная расчетная область $\Omega$ (рис. 1) описывается как объединение двух прямоугольников: основания и <<шапки>>.

\begin{figure}[H]
    \centering
    \includegraphics[width=0.6\textwidth]{image_d6a054.png}
    \caption{Расчетная Т-образная область}
    \label{fig:domain}
\end{figure}

Допускается использовать фиктивные узлы для сохранения регулярной структуры СЛАУ. Фиктивные узлы лежат вне Т-образной области. Для таких узлов матричное уравнение вырождается в тривиальное $u_{i,j} = 0$.

\subsection{Сборка СЛАУ и метод решения}
Для хранения глобальной матрицы $A$ используется диагональный формат (DIA) с пятью ненулевыми диагоналями.
Для решения полученной СЛАУ применяется итерационный метод блочной релаксации (метод прогонки по линиям). Итерационный процесс организован по горизонтальным линиям сетки. Для каждой строки $j$ формируется и решается методом прогонки трехдиагональная подсистема:
\begin{equation}
a_i u_{i-1,j}^{(k+1)} + c_i u_{i,j}^{(k+1)} + b_i u_{i+1,j}^{(k+1)} = F(u_{i,j-1}^{(k+1)}, u_{i,j+1}^{(k)})
\end{equation}

\section{Тестирование и порядок аппроксимации}

Тестирование программы проводилось на аналитических решениях в виде полиномов первой и второй степени. Результаты тестирования представлены в таблице \ref{tab:poly}.

\begin{table}[H]
    \centering
    \caption{Тестирование программы на полиномах ($\lambda=1, \gamma=1, \beta=2$)}
    \label{tab:poly}
    \begin{tabularx}{\textwidth}{|c|c|c|X|}
        \hline
        \textbf{Точное решение $u^*$} & \textbf{Шаг сетки $h$} & \textbf{Норма ошибки $\|u^* - u\|_\infty$} & \textbf{Комментарий} \\
        \hline
        $u(x,y) = 2x + 3y + 1$ & 0.1 & [ВСТАВИТЬ $10^{-15}$] & Схема точна для линейных полиномов. \\
        \hline
        $u(x,y) = x^2 + y^2$ & 0.1 & [ВСТАВИТЬ ОШИБКУ] & Появление ошибки из-за 1-го порядка аппроксимации краевых условий 3-го рода. \\
        \hline
    \end{tabularx}
\end{table}

Исследование порядка аппроксимации реализованного метода проводилось на неполиномиальном решении $u(x,y) = \sin(x+y)$. Экспериментальный порядок сходимости $p$ вычисляется по формуле: $p = \log_2 (E_h / E_{h/2})$.

\begin{table}[H]
    \centering
    \caption{Исследование порядка аппроксимации}
    \label{tab:order}
    \begin{tabular}{|c|c|c|c|}
        \hline
        \textbf{Размер сетки $N_x \times N_y$} & \textbf{Шаг $h$} & \textbf{Ошибка $E_h$} & \textbf{Порядок $p$} \\
        \hline
        $10 \times 10$ & $h_0$ & [ВСТАВИТЬ] & - \\
        \hline
        $20 \times 20$ & $h_0 / 2$ & [ВСТАВИТЬ] & [ВСТАВИТЬ] \\
        \hline
        $40 \times 40$ & $h_0 / 4$ & [ВСТАВИТЬ] & [ВСТАВИТЬ] \\
        \hline
    \end{tabular}
\end{table}

\section{Текст программы}

\begin{lstlisting}
#include <iostream>
#include <vector>
#include <cmath>
#include <iomanip>
#include <algorithm>

using namespace std;

// Параметры задачи
struct Params {
    double lambda = 1.0;
    double gamma = 1.0;
    double beta = 2.0;
};

// Типы узлов
enum NodeType { FICTITIOUS = 0, INTERNAL = 1, BOUNDARY_1 = 2, BOUNDARY_3 = 3 };

// Точное решение и правая часть
double u_exact(double x, double y) { return sin(x + y); }
double f_func(double x, double y, Params p) { 
    return p.lambda * 2.0 * sin(x + y) + p.gamma * u_exact(x, y); 
}

// Метод прогонки
void tridiagonal_solve(const vector<double>& a, const vector<double>& c, 
                       const vector<double>& b, const vector<double>& f, vector<double>& x) {
    int n = f.size();
    vector<double> alpha(n, 0.0), beta(n, 0.0);
    
    alpha[0] = -b[0] / c[0];
    beta[0] = f[0] / c[0];
    
    for (int i = 1; i < n - 1; ++i) {
        double denom = c[i] + a[i] * alpha[i - 1];
        alpha[i] = -b[i] / denom;
        beta[i] = (f[i] - a[i] * beta[i - 1]) / denom;
    }
    
    x[n - 1] = (f[n - 1] - a[n - 1] * beta[n - 2]) / (c[n - 1] + a[n - 1] * alpha[n - 2]);
    for (int i = n - 2; i >= 0; --i) {
        x[i] = alpha[i] * x[i + 1] + beta[i];
    }
}

// Блочная релаксация
void block_relaxation(const vector<double>& diMain, const vector<double>& diLower, 
                      const vector<double>& diUpper, const vector<double>& diFarLower, 
                      const vector<double>& diFarUpper, const vector<double>& b, 
                      vector<double>& u, int Nx, int Ny, double eps = 1e-9) {
    int N = Nx * Ny;
    vector<double> u_new = u;
    double max_diff;
    int iter = 0;
    
    do {
        max_diff = 0.0;
        for (int j = 0; j < Ny; ++j) {
            vector<double> a(Nx, 0.0), c(Nx, 0.0), b_sub(Nx, 0.0), f(Nx, 0.0), x(Nx, 0.0);
            for (int i = 0; i < Nx; ++i) {
                int k = j * Nx + i;
                c[i] = diMain[k];
                if (i > 0) a[i] = diLower[k];
                if (i < Nx - 1) b_sub[i] = diUpper[k];
                
                f[i] = b[k];
                if (j > 0) f[i] -= diFarLower[k] * u_new[k - Nx];
                if (j < Ny - 1) f[i] -= diFarUpper[k] * u[k + Nx];
            }
            tridiagonal_solve(a, c, b_sub, f, x);
            for (int i = 0; i < Nx; ++i) {
                int k = j * Nx + i;
                max_diff = max(max_diff, abs(x[i] - u[k]));
                u_new[k] = x[i];
            }
        }
        u = u_new;
        iter++;
    } while (max_diff > eps && iter < 10000);
}

int main() {
    Params p;
    // Геометрия Т-области
    double cap_x_min = 0.0, cap_x_max = 6.0, cap_y_min = 3.0, cap_y_max = 5.0;
    double base_x_min = 2.0, base_x_max = 4.0, base_y_min = 0.0, base_y_max = 3.0;

    int Nx = 40, Ny = 40;
    double hx = 6.0 / (Nx - 1), hy = 5.0 / (Ny - 1);
    int N = Nx * Ny;
    
    vector<NodeType> node_types(N, FICTITIOUS);
    
    // Разметка узлов
    for (int j = 0; j < Ny; ++j) {
        for (int i = 0; i < Nx; ++i) {
            double x = i * hx;
            double y = j * hy;
            int k = j * Nx + i;
            
            bool in_cap = (x >= cap_x_min - 1e-7 && x <= cap_x_max + 1e-7 && y >= cap_y_min - 1e-7 && y <= cap_y_max + 1e-7);
            bool in_base = (x >= base_x_min - 1e-7 && x <= base_x_max + 1e-7 && y >= base_y_min - 1e-7 && y <= base_y_max + 1e-7);

            if (in_cap || in_base) {
                if (abs(y - cap_y_max) < 1e-7 || abs(y - base_y_min) < 1e-7) {
                    node_types[k] = BOUNDARY_3; // 3-й род верх и низ
                } else if ((abs(x - cap_x_min) < 1e-7 && y >= cap_y_min) || (abs(x - cap_x_max) < 1e-7 && y >= cap_y_min) ||
                           (abs(x - base_x_min) < 1e-7 && y <= base_y_max) || (abs(x - base_x_max) < 1e-7 && y <= base_y_max) ||
                           (abs(y - cap_y_min) < 1e-7 && (x <= base_x_min + 1e-7 || x >= base_x_max - 1e-7))) {
                    node_types[k] = BOUNDARY_1;
                } else {
                    node_types[k] = INTERNAL;
                }
            }
        }
    }

    // Сборка СЛАУ
    vector<double> diMain(N, 0.0), diLower(N, 0.0), diUpper(N, 0.0), diFarLower(N, 0.0), diFarUpper(N, 0.0), b(N, 0.0);
    
    for (int j = 0; j < Ny; ++j) {
        for (int i = 0; i < Nx; ++i) {
            int k = j * Nx + i;
            double x = i * hx, y = j * hy;

            if (node_types[k] == FICTITIOUS) {
                diMain[k] = 1.0; b[k] = 0.0;
            } else if (node_types[k] == BOUNDARY_1) {
                diMain[k] = 1.0; b[k] = u_exact(x, y);
            } else if (node_types[k] == BOUNDARY_3) {
                if (abs(y - base_y_min) < 1e-7) { // Низ
                    diMain[k] = p.lambda / hy + p.beta;
                    diFarUpper[k] = -p.lambda / hy;
                    b[k] = p.beta * u_exact(x, y); // Упрощенно u_beta = u_exact
                } else { // Верх
                    diMain[k] = p.lambda / hy + p.beta;
                    diFarLower[k] = -p.lambda / hy;
                    b[k] = p.beta * u_exact(x, y);
                }
            } else if (node_types[k] == INTERNAL) {
                diMain[k] = -(2.0 * p.lambda / (hx * hx) + 2.0 * p.lambda / (hy * hy)) + p.gamma;
                diLower[k] = p.lambda / (hx * hx);
                diUpper[k] = p.lambda / (hx * hx);
                diFarLower[k] = p.lambda / (hy * hy);
                diFarUpper[k] = p.lambda / (hy * hy);
                b[k] = -f_func(x, y, p);
            }
        }
    }

    vector<double> u(N, 0.0);
    block_relaxation(diMain, diLower, diUpper, diFarLower, diFarUpper, b, u, Nx, Ny);

    // Подсчет ошибки
    double max_err = 0.0;
    for (int j = 0; j < Ny; ++j) {
        for (int i = 0; i < Nx; ++i) {
            int k = j * Nx + i;
            if (node_types[k] != FICTITIOUS) {
                max_err = max(max_err, abs(u[k] - u_exact(i * hx, j * hy)));
            }
        }
    }
    cout << "Grid: " << Nx << "x" << Ny << ", Max Error: " << max_err << endl;

    return 0;
}
\end{lstlisting}

\section{Выводы}
В ходе выполнения лабораторной работы была реализована программа на языке С++ для решения двумерной эллиптической краевой задачи методом конечных разностей. 
\begin{enumerate}
    \item Использование фиктивных узлов позволило сохранить регулярную структуру конечно-разностной сетки для непрямоугольной расчетной области и применить метод блочной релаксации (прогонки по линиям) для эффективного решения полученной СЛАУ.
    \item Результаты тестирования показали, что при наличии краевых условий третьего рода, аппроксимируемых односторонними разностями первого порядка, глобальный порядок сходимости всей схемы падает до первого порядка ($O(h)$), несмотря на то, что аппроксимация во внутренних узлах имеет второй порядок.
\end{enumerate}

\end{document}